\section{Introduction}

The computational cost of training state-of-the-art foundation models has been increasing rapidly~\citep{amodei2018ai,sevilla2022compute}. Since pretraining modern Transformer-based models is dominated by dense matrix multiplications, reducing the precision of these operations is one of the most direct levers for improving training efficiency. This has motivated a progression from FP16/BF16 to FP8 training \citep{micikevicius2022fp8}, and, more recently, to hardware-supported 4-bit block scale formats \citep{rouhani2023microscalingdataformatsdeep} such as MXFP4 and NVFP4 on NVIDIA Blackwell GPUs \citep{nvidia2025blackwell, alvarez2025introducingnvfp4, nvidia2025nvfp4}.

Yet, quantized training introduces noise, and in the 4-bit regime this noise could affect convergence, even when operations may appear accurate. A recent line of work has therefore focused on recipes that recover most of the accuracy of FP8 or BF16 training while preserving the speedups of 4-bit tensor-core execution. Examples include the \nvrecipe recipe \citep{nvidia2025pretraining}, backward-pass FP4 training \citep{tseng2025mxfp4}, fully quantized FP4 training \citep{chmiel2025fp4all}, TetraJet \citep{chen2025oscillationreducedmxfp4trainingvision}, TetraJet-v2 \citep{chen2025tetrajet}, Four Over Six \citep{cook2025fouroversix}, Quartet \citep{castro2025quartet}, and Quartet II \citep{panferov2026quartetii}. Despite this rapid progress, it is fair to say that there is still no definitive solution: low-precision training can still induce small but persistent accuracy degradation, and existing methods disagree about which techniques are essential.

At the same time, some principles are starting to emerge. On the forward pass, multiple methods implicitly or explicitly optimize representation quality, for instance by substituting the \nvrecipe{} recipe's two-dimensional weight quantization \citep{nvidia2025pretraining} with one-dimensional quantization for finer-grained scaling or by reducing quantization error for weights and activations through improved scale selection \citep{chen2025oscillationreducedmxfp4trainingvision, castro2025quartet, chen2025tetrajet, cook2025fouroversix, panferov2026quartetii}. On the backward pass, a line of evidence suggests that bias in the gradient estimator can accumulate over long training horizons, making unbiased or carefully controlled quantization important for stable convergence \citep{chmiel2024accurateneuraltraining4bit, tseng2025mxfp4, chen2025oscillationreducedmxfp4trainingvision, chmiel2025fp4all, nvidia2025pretraining, chen2025tetrajet, panferov2026quartetii}. The difficulty is that these principles are not always clearly compatible. For instance, the \nvrecipe recipe uses two-dimensional quantization of weights to maintain consistency between forward and backward quantized weight representations \citep{nvidia2025pretraining}. By contrast, Quartet II introduces a technique to address bias in the backward pass that relies on random rotations applied along the inner dimension of both backward GEMMs; those rotations, in turn, are most naturally paired with one-dimensional quantization \citep{panferov2026quartetii, tseng2025mxfp4, chen2025tetrajet}. To approximate consistency with the forward-pass weight representation, one-dimensional schemes often dequantize the forward-quantized weights and re-quantize them along the backward contraction dimension rather than directly quantizing the high-precision weights during the backward pass \citep{chen2025oscillationreducedmxfp4trainingvision, castro2025quartet, chen2025tetrajet, panferov2026quartetii}.

This creates a practical question: which ingredients are non-negotiable, and which ones are merely useful at a given scale? The key problem we aim to address here is that the ``right'' answer is hard to infer from short ablations. Many accuracy effects in low-precision training only become visible late in training, or only after increasing model scale and token budget. As a result, as we show, recipes can appear close early on while diverging later, in ways that are expensive to diagnose. 

In this paper, we aim to isolate a set of  principles that matter for NVFP4 training by systematically comparing variants of existing recipes at medium scale. We evaluate dense 8B and 30B-A3B mixture-of-experts (MoE) models from the Nemotron  model family, trained up to 1T tokens. We organize our comparison around the major design axes that separate current methods: whether to use selective high-precision layers, how to select scales in the forward and backward pass, whether to quantize weights in one or two dimensions, strategies to mitigate bias in the backward pass, and how these choices interact.

\paragraph{Contributions.}
Our main finding is that the current recipes capture complementary pieces of the problem, and we propose a ``Mosaic'' approach that blends the best qualities of each. In more detail:
\begin{enumerate}
    \item We perform the first detailed and unified ablation of NVFP4 training techniques, examining how each technique reduces the model-quality gap relative to BF16 training.  
    \item  
    We show that keeping selected layers in higher precision is necessary for    
  full accuracy recovery in NVFP4 training, even with an optimized recipe. Omitting this can cause training instabilities at larger model sizes and medium-scale token horizons.    
    \item 
     We find that selectively applying the MS-EDEN \citep{panferov2026quartetii} unbiased estimator to gradients and Four Over Six \citep{cook2025fouroversix} scale selection to the remaining backward tensors yields the lowest observed loss gap among the tested backward-pass configurations when all layers are quantized.
    \item 
    Based on these observations, we develop a combined NVFP4 training recipe, which we call ``Mosaic NVFP4,'' with selective high-precision layers, Random Hadamard Transforms on both backward GEMMs, unbiased gradients via MS-EDEN, proposed by Quartet II, and Four Over Six scale selection on remaining forward and backward tensors. To the best of our knowledge, this is the strongest FP4 training result demonstrated at this scale in loss gap to BF16. This scheme remains compatible with efficient kernels. 
\end{enumerate}

